\documentclass[11pt]{article}
\usepackage[a4paper, margin=1in]{geometry}
\usepackage[T1]{fontenc}
\usepackage[utf8]{inputenc}
\usepackage{amsmath, amssymb}
\usepackage{microtype}
\usepackage{parskip}
\usepackage{hyperref}
\hypersetup{colorlinks=true, urlcolor=blue}
\pagestyle{empty}

\begin{document}

\begin{center}
{\Large\bfseries Quantum Mechanics without Ontology?}\\[0.4em]
{\large Configuration Space, the Arena Problem, and the Road Not Taken in 1926}
\end{center}

\vspace{1.5em}

\noindent\textbf{Author.} Claes Johnson\\
Professor emeritus of Applied Mathematics\\
KTH Royal Institute of Technology, SE-100 44 Stockholm, Sweden\\
\texttt{claesjohnson@gmail.com}\\
ORCID / affiliation details as required by the journal.

\vspace{1em}

\noindent\textbf{Corresponding author.} Claes Johnson (\texttt{claesjohnson@gmail.com}).

\vspace{1em}

\noindent\textbf{Declarations.}
\begin{itemize}
\item \emph{Funding:} none.
\item \emph{Competing interests:} the author is the proponent of the three-dimensional charge-density
formulation (RealQM) discussed in the paper as one instance of a primitive-ontology approach.
\item \emph{Use of AI.} The manuscript was prepared in cooperation with an AI assistant (Claude, Anthropic),
which contributed to drafting and to surveying the secondary literature. The argument, the claims, and the
responsibility for the content are the author's. The author is glad to describe this collaboration further if
the journal wishes.
\end{itemize}

\vspace{1em}

\noindent\textbf{Abstract.}
It is often said that quantum mechanics dispenses with ontology and that its predictive success licenses the
omission. This paper argues that the situation is both more precise and more serious. Under the
preparation-independence assumption, the Pusey--Barrett--Rudolph theorem tells against a purely epistemic reading
of the quantum state, so the state is naturally taken to be ontic; the genuine difficulty is not the existence of
an ontology but its \emph{arena} --- the state is a function on $3N$-dimensional configuration space, not on the
three-dimensional space that observers and instruments occupy. The measurement problem, the status of collapse,
and the proliferation of interpretations are read as symptoms of this arena problem. The paper revisits the
historical moment (spring 1926) at which the wave function was lifted from three-dimensional space to
configuration space over the objections of Schr\"odinger and Einstein; situates the options within the
contemporary debate between wave-function realism and primitive ontology; presents a three-dimensional
charge-density formulation as an under-explored primitive-ontology alternative --- a concrete realization of the
road not taken; and draws a practical corollary from quantum computing, whose claimed advantage turns on the
physical reality of the exponentially large amplitude.

\vspace{1em}

\noindent\textbf{Keywords.} quantum state; ontology; configuration space; wave-function realism; primitive
ontology; PBR theorem; measurement problem; history of quantum mechanics.

\end{document}
